\section{Arithmetic orbit groupoids and toric period morphisms}
\label{sec:groupoid-motivic-interfaces}

This section constructs two interfaces which begin with the same
positive-odd path coordinates but retain different information.  The first
is the discrete orbit groupoid of the actual positive-integer map.  Its
singletons are compact open bisections, so the finite-support path algebra
embeds by point masses.  The second is the germ groupoid generated by exact
partial homeomorphisms of $\mathbb Z_3$.  Its unit space is compact and its
arrows vary $3$-adically; arithmetic paths enter it by a faithful functor,
not by identifying its convolution algebra with the discrete path algebra.

The period interface is equally explicit.  A unimodular change of torus
coordinates identifies the path bidegree $(\ell,A-\ell)$ with the character
$(A,-\ell)$, and every logarithm used below is obtained from a commutative
square of Deligne $1$-motives.  The arctangent integral is related to a
Kummer logarithm by an isomorphism of algebraic pairs.  These are actual
maps; no categorical relation is inferred from similar formulas.

Deaconu's construction is the literature prototype for a groupoid attached
to an endomorphism, but his topological theorem assumes a compact Hausdorff
space and a self-covering map \cite[printed pp.~1779--1781]{Deaconu1995}.
Our discrete theorem is therefore proved directly.  The definitions and
period computations for $1$-motives are crosswalked against
Huber--W\"ustholz \cite[Definitions~8.1 and 10.2; Proposition~12.10]{HuberWustholz2022};
the permitted arrows in Nori's diagram of pairs are those of
Huber--M\"uller-Stach \cite[Definitions~7.1.4 and 9.1.1]{HuberMullerStach2025}.
Barbieri-Viale--Kahn supply the general derived and Cartier-dual framework
\cite[Section~1.1 and Theorem~2.1.2]{BarbieriVialeKahn2016}.  None of these
sources states the Collatz-specific constructions proved here.

\subsection{The discrete arithmetic orbit groupoid}

Retain the positive-odd map $C$, exponent function $a$, and first-return
category $\mathcal E_C$ from
Definition~\ref{def:odd-first-return-path-category}.  For $n\geq0$ put
\[
 S_na(x)=\sum_{j=0}^{n-1}a(C^jx),
 \qquad S_0a(x)=0.
\]

\begin{definition}[arithmetic orbit groupoid]
\label{def:arithmetic-orbit-groupoid}
Let
\begin{equation}
\label{eq:arithmetic-orbit-groupoid}
 \mathcal G_{\mathrm{arith}}
 =\left\{(x,k,y)\in\mathcal O\times\mathbb Z\times\mathcal O:
 \begin{array}{l}
 \text{there are }m,n\geq0\text{ with}\ k=m-n,\\[-1mm]
 C^m(x)=C^n(y)
 \end{array}\right\}.
\end{equation}
Its range and source are $r(x,k,y)=x$ and $s(x,k,y)=y$, and
\[
 (x,k,y)(y,l,z)=(x,k+l,z),
 \qquad (x,k,y)^{-1}=(y,-k,x).
\]
We give $\mathcal G_{\mathrm{arith}}$ the discrete topology.
\end{definition}

\begin{theorem}[exact arithmetic groupoid and its two cocycles]
\label{thm:arithmetic-groupoid-cocycles}
The preceding formulas define a countable, Hausdorff, locally compact,
\'{e}tale groupoid.  Counting measure on every range fibre is a Haar system.
For $g=(x,k,y)$ choose any $m,n\geq0$ representing it and put
\begin{equation}
\label{eq:arithmetic-groupoid-ca}
 c_a(g)=S_ma(x)-S_na(y).
\end{equation}
This integer does not depend on the representative.  The maps
\begin{align}
 \kappa(g)&=k,\label{eq:arithmetic-groupoid-kappa}\\
 D(g)&=(-k,-c_a(g)+k)\in\mathbb Z^2,\label{eq:arithmetic-groupoid-D}\\
 \delta_{\mathcal G}(g)&=-c_a(g)\log2+k\log3\label{eq:arithmetic-groupoid-delta}
\end{align}
are groupoid homomorphisms to the displayed additive groups.
\end{theorem}

\begin{proof}
If $(m,n)$ represents $(x,k,y)$ and $t\geq0$, then
$(m+t,n+t)$ represents the same triple.  With
$u=C^m(x)=C^n(y)$, the Birkhoff-sum identity gives
\[
 S_{m+t}a(x)=S_ma(x)+S_ta(u),\qquad
 S_{n+t}a(y)=S_na(y)+S_ta(u),
\]
so their difference is unchanged.  Any two representatives have the same
difference $m-n$; after extending the shorter pair, both reach a common pair
of times.  This proves that \eqref{eq:arithmetic-groupoid-ca} is
well-defined.

For multiplication, suppose $(m,n)$ represents $g=(x,k,y)$ and $(p,q)$
represents $h=(y,l,z)$.  Put $L=\max\{n,p\}$.  Extending the first pair by
$L-n$ and the second by $L-p$ gives a representative
\[
 \bigl(m+L-n,\ q+L-p\bigr)
\]
of $gh$.  The two sums along the common middle orbit cancel, giving
$c_a(gh)=c_a(g)+c_a(h)$.  The same aligned representative gives
$\kappa(gh)=\kappa(g)+\kappa(h)$.  Equations
\eqref{eq:arithmetic-groupoid-D} and
\eqref{eq:arithmetic-groupoid-delta} are integral and real linear
combinations of these two cocycles, so they are additive as well.

The groupoid axioms now follow from the displayed product and inverse.
The underlying set is countable.  In the discrete topology every singleton
is a compact open bisection, range and source are local homeomorphisms, and
compact subsets are finite.  Thus the groupoid is Hausdorff, locally compact,
and \'{e}tale.  Left multiplication bijects range fibres, so counting measures
are left invariant and form the counting Haar system.
\end{proof}

\begin{theorem}[faithful path embedding and operator actions]
\label{thm:path-to-arithmetic-groupoid}
For a path $p:x\to y$ define
\begin{equation}
\label{eq:path-groupoid-embedding}
 J(p)=(y,-\ell(p),x).
\end{equation}
Then $J\colon\mathcal E_C\to\mathcal G_{\mathrm{arith}}$ is an injective
functor, where the target is regarded as a groupoid category.  On its image,
\begin{equation}
\label{eq:path-groupoid-cocycle-recovery}
 D(J(p))=(\ell(p),A(p)-\ell(p)),
 \qquad
 \delta_{\mathcal G}(J(p))=A(p)\log2-\ell(p)\log3.
\end{equation}
Consequently the point-mass map
\[
 \mathbb C[\mathcal E_C]\longrightarrow C_c(\mathcal G_{\mathrm{arith}}),
 \qquad [p]\longmapsto 1_{\{J(p)\}},
\]
is an injective algebra homomorphism in the target-left convention.

For $z=(z_1,z_2)\in\mathbb T^2$ and $t\in\mathbb R$, the formulas
\begin{align}
 (\Gamma_zf)(g)&=z_1^{D_1(g)}z_2^{D_2(g)}f(g),
 \label{eq:arithmetic-groupoid-gauge}\\
 (\alpha_tf)(g)&=e^{it\delta_{\mathcal G}(g)}f(g)
 \label{eq:arithmetic-groupoid-time}
\end{align}
define commuting automorphism groups of the convolution algebra.  They
extend separately to both the full and reduced groupoid $C^*$-algebras.
\end{theorem}

\begin{proof}
The pair $(0,\ell(p))$ represents $J(p)$ because
$C^0(y)=C^{\ell(p)}(x)$.  If $p:x\to y$ and $q:z\to x$, then
\[
 J(p)J(q)=(y,-\ell(p)-\ell(q),z)=J(p\circ q),
\]
which is exactly the target-left multiplication used in
\eqref{eq:weighted-path-target-left-product}.  A path of the deterministic
category is fixed by its source and length, so $J$ is injective.  The same
representative has $c_a(J(p))=-A(p)$; substitution gives
\eqref{eq:path-groupoid-cocycle-recovery}.  Distinct point masses are
linearly independent, proving the algebra assertion.

Additivity of $D$ and $\delta_{\mathcal G}$ makes
\eqref{eq:arithmetic-groupoid-gauge} and
\eqref{eq:arithmetic-groupoid-time} multiplicative under convolution and
compatible with involution.  Their inverses are obtained from $z^{-1}$ and
$-t$.  The universal property gives the full extensions.  In every regular
representation the same formulas are implemented by diagonal unitaries, so
they preserve the reduced norm and give the reduced extensions.
\end{proof}

\begin{nonclaim}
No amenability assertion or equality between the full and reduced
$C^*$-algebras is used.  Equations
\eqref{eq:arithmetic-groupoid-gauge}--\eqref{eq:arithmetic-groupoid-time}
construct dynamics; they do not classify equilibrium states or prove a phase
transition.
\end{nonclaim}

\subsection{The compact 3-adic pseudogroup and the arithmetic functor}

Set $X=\mathbb Z_3$.  For $a\geq1$ put
\begin{equation}
\label{eq:three-adic-branch-pair}
 Y_a=\{y\in X:2^ay\equiv1\pmod3\},\qquad
 g_a(x)=\frac{3x+1}{2^a},\qquad
 \beta_a(y)=\frac{2^ay-1}{3}.
\end{equation}

\begin{proposition}[exact $3$-adic branch domains]
\label{prop:three-adic-branch-homeomorphisms}
The set $Y_a$ is the clopen residue class
$2^{-a}+3\mathbb Z_3$.  The maps
\[
 g_a:X\longrightarrow Y_a,
 \qquad \beta_a:Y_a\longrightarrow X
\]
are mutually inverse homeomorphisms.  For a chronological word
$w=(a_1,\ldots,a_n)$, the composite
\[
 h_w=g_{a_n}\circ\cdots\circ g_{a_1}:X\longrightarrow Y_w:=h_w(X)
\]
is a homeomorphism onto a clopen set and has the affine formula
\begin{equation}
\label{eq:three-adic-word-affine}
 h_w(x)=\frac{3^nx+B_w}{2^{A(w)}},
 \qquad
 B_w=\sum_{i=1}^n3^{n-i}2^{a_1+\cdots+a_{i-1}}.
\end{equation}
\end{proposition}

\begin{proof}
The congruence defining $Y_a$ is one residue class because $2^a$ is a unit
in $\mathbb Z_3$.  The numerator of $\beta_a$ lies in $3\mathbb Z_3$ on
that class, and direct substitution gives
\[
 \beta_ag_a(x)=x,\qquad g_a\beta_a(y)=y.
\]
Both maps are continuous affine maps, hence are inverse homeomorphisms.
Composition proves the topological assertion for words; induction gives
\eqref{eq:three-adic-word-affine}.
\end{proof}

Let $\mathcal P_3$ be the pseudogroup generated by restrictions of the
$g_a$ and $\beta_a$, and let $\mathcal H_3$ be its groupoid of germs.  We
write $[\phi,x]$ for the germ at $x$ of a pseudogroup element $\phi$.

\begin{theorem}[the $3$-adic germ groupoid and its arithmetic subgroupoid]
\label{thm:three-adic-germ-functor}
The groupoid $\mathcal H_3$ is second countable, Hausdorff, locally compact,
and \'{e}tale, with compact unit space $X$ and its counting Haar system.

For $g=(x,m-n,y)\in\mathcal G_{\mathrm{arith}}$, let $p_m(x)$ and $p_n(y)$
be the actual chronological exponent words of the first $m$ and $n$ steps.
Then
\begin{equation}
\label{eq:arithmetic-to-germ-functor}
 \Xi(g)=
 \left[h_{p_m(x)}^{-1}\circ h_{p_n(y)},y\right]
\end{equation}
is a well-defined injective groupoid functor, with object map the inclusion
$\mathcal O\hookrightarrow\mathbb Z_3$,
$\Xi\colon\mathcal G_{\mathrm{arith}}\to\mathcal H_3$.  It sends
$J(p)$ to $[h_p,s(p)]$.

If $u(g)$ is the rational slope of an affine representative of $\Xi(g)$,
then
\begin{equation}
\label{eq:germ-slope-cocycles}
 u(g)=3^{n-m}2^{S_ma(x)-S_na(y)},\qquad
 D(g)=\bigl(\nu_3(u(g)),-\nu_2(u(g))-\nu_3(u(g))\bigr),
\end{equation}
and
\begin{equation}
\label{eq:germ-real-cocycle}
 \delta_{\mathcal G}(g)=-\log|u(g)|_{\infty}.
\end{equation}
Thus $\Xi$ preserves both the bidegree and real cocycle.
\end{theorem}

\begin{proof}
Every element of $\mathcal P_3$ is locally the restriction of an affine map
$x\mapsto ux+v$ with $u\ne0$.  Two such maps agreeing on a nonempty
$3$-adic open set have the same slope and intercept.  Hence two distinct
germs with a common source can be separated by basic germ bisections; germs
with different sources are separated in the unit space.  This proves
Hausdorffness.  The usual germ bisections make range and source local
homeomorphisms.  Since $\mathbb Z_3$ has a countable clopen basis and the
generating words are countable, the groupoid is second countable.  Every
germ has a bisection with compact-open source, proving local compactness.
Counting measures are the canonical Haar system of an \'{e}tale groupoid.

The equality $C^m(x)=C^n(y)$ says that the two word maps in
\eqref{eq:arithmetic-to-germ-functor} meet at their displayed points, so the
composite germ is defined near $y$ and maps it to $x$.  Replacing $(m,n)$ by
$(m+t,n+t)$ prefixes both word maps by the same future affine map $h_r$;
the factor $h_r$ cancels between inverse and forward maps.  Two arbitrary
representatives admit a common extension, proving well-definedness.  The
same common-middle alignment used in
Theorem~\ref{thm:arithmetic-groupoid-cocycles} proves functoriality.

The slope of $h_{p_m(x)}$ is $3^m2^{-S_ma(x)}$, so the slope of its inverse
followed by $h_{p_n(y)}$ is the first expression in
\eqref{eq:germ-slope-cocycles}.  Its $3$-adic valuation is $n-m$, which
recovers the integer coordinate $m-n$.  Source, range, and this coordinate
recover the entire triple $(x,m-n,y)$, so $\Xi$ is injective.  The two
valuation formulas and \eqref{eq:germ-real-cocycle} follow by substitution.
For $J(p)$, the representative is $(0,\ell(p))$, so its germ is $h_p$.
\end{proof}

The functor $\Xi$ is not an inclusion of convolution algebras by point
masses: a singleton germ over the nondiscrete unit space $\mathbb Z_3$ is
generally not open, so its characteristic function is not in
$C_c(\mathcal H_3)$.  It is also not legitimate to obtain an operator model
on $\ell^2(\mathcal O)$ by restricting every $g_a$.

\begin{proposition}[explicit failure of the unrestricted arithmetic
restriction]
\label{prop:three-adic-l2-restriction-failure}
The $3$-adic homeomorphism $g_a:X\to Y_a$ does not preserve
$\mathcal O$.  In particular $g_1(1)=2$.  If one defines
\[
 (V_a\xi)(y)=
 \begin{cases}
  \xi(\beta_a(y)),&y\in Y_a\cap\mathcal O,\\
  0,&y\notin Y_a\cap\mathcal O,
 \end{cases}
\]
on $\ell^2(\mathcal O)$, then $V_1\delta_1=0$ and hence
$V_1^*V_1\ne1$.
\end{proposition}

\begin{proof}
The equation $\beta_1(y)=1$ has the unique solution $y=2$, which is not a
positive odd integer.  Therefore every coordinate of $V_1\delta_1$ vanishes.
\end{proof}

There is a second elementary obstruction to a commonly proposed operator
symmetry.  For a nonempty free word $q$, left prefixing
$L_q[p]=[qp]$ is not an algebra endomorphism of the free-monoid
algebra, since
\[
 L_q(1\cdot1)=[q]\ne[qq]=L_q(1)L_q(1).
\]
Consequently no prefix endomorphism of a completion follows from word
concatenation alone; its action on the diagonal and all covariance relations
would have to be declared and proved separately.

\begin{nonclaim}
The exact objects constructed here are $\mathcal G_{\mathrm{arith}}$,
$\mathcal H_3$, and the faithful functor $\Xi$.  No unspecified semigroup
crossed product is identified with either full or reduced groupoid
$C^*$-algebra.  Such an identification requires a declared covariance
relation and a proof of the corresponding universal property.
\end{nonclaim}

\subsection{The exact torus coordinate change}

Work over $\mathbb Q$.  Use two copies of the split torus, with coordinates
\[
 T_{\mathrm{ar}}=(\mathbf G_m)^2_{x,y},\qquad
 T_{\deg}=(\mathbf G_m)^2_{\lambda,\mu}.
\]
For a path $p$ put
\[
 \chi_p(x,y)=x^{A(p)}y^{-\ell(p)},\qquad
 \psi_p(\lambda,\mu)=\lambda^{\ell(p)}\mu^{A(p)-\ell(p)}.
\]

\begin{theorem}[unimodular arithmetic-degree torus isomorphism]
\label{thm:torus-coordinate-isomorphism}
The maps
\begin{equation}
\label{eq:torus-coordinate-isomorphism}
 \Phi:T_{\mathrm{ar}}\longrightarrow T_{\deg},
 \quad \Phi(x,y)=(x/y,x),
 \qquad
 \Phi^{-1}(\lambda,\mu)=(\mu,\mu/\lambda)
\end{equation}
are inverse torus isomorphisms, and
\begin{equation}
\label{eq:torus-character-intertwining}
 \psi_p\circ\Phi=\chi_p.
\end{equation}
On exponent lattices the mutually inverse matrices are
\begin{equation}
\label{eq:torus-exponent-matrices}
 \binom{A}{-\ell}
 =\begin{pmatrix}1&1\\ -1&0\end{pmatrix}
  \binom{\ell}{A-\ell},
 \qquad
 \binom{\ell}{A-\ell}
 =\begin{pmatrix}0&-1\\ 1&1\end{pmatrix}
  \binom{A}{-\ell}.
\end{equation}
\end{theorem}

\begin{proof}
Direct substitution proves the inverse formulas and
\[
 \psi_p(\Phi(x,y))=(x/y)^{\ell(p)}x^{A(p)-\ell(p)}
 =x^{A(p)}y^{-\ell(p)}.
\]
Multiplication of the two integral matrices gives the identity in both
orders, proving the lattice statement without changing either coordinate
system.
\end{proof}

Define the two isomorphic $1$-motives
\begin{align}
 M_{\mathrm{ar}}&=
 [\mathbb Z\xrightarrow{u_{\mathrm{ar}}}T_{\mathrm{ar}}],
 &u_{\mathrm{ar}}(n)&=(2^n,3^n),\label{eq:arithmetic-one-motive}\\
 M_{\deg}&=
 [\mathbb Z\xrightarrow{u_{\deg}}T_{\deg}],
 &u_{\deg}(n)&=((2/3)^n,2^n).\label{eq:degree-one-motive}
\end{align}
Then $(\mathrm{id}_{\mathbb Z},\Phi)$ is an isomorphism
$M_{\mathrm{ar}}\to M_{\deg}$.

\begin{theorem}[character-induced Kummer morphisms and exact periods]
\label{thm:kummer-character-morphisms}
For a path $p$ set
\[
 \alpha_p=2^{A(p)}3^{-\ell(p)},
 \qquad
 K(\alpha_p)=[\mathbb Z\xrightarrow{n\mapsto\alpha_p^n}\mathbf G_m].
\]
Then
\begin{equation}
\label{eq:kummer-character-morphism}
 (\mathrm{id}_{\mathbb Z},\chi_p):M_{\mathrm{ar}}\longrightarrow K(\alpha_p)
\end{equation}
is a morphism of $1$-motives.  Along the positive-real path
$\gamma(s)=(2^s,3^s)$, $0\leq s\leq1$,
\begin{equation}
\label{eq:kummer-path-period}
 \int_\gamma\chi_p^*\frac{dz}{z}
 =A(p)\log2-\ell(p)\log3.
\end{equation}
The branch is fixed by this path.

For every prefix $1\leq k\leq\ell(p)$, the character
\[
 \chi_{p,k}(x,y)=x^{-A_k(p)}y^{k-1}
\]
similarly gives a morphism to $K(r_k(p))$, where
$r_k(p)=3^{k-1}2^{-A_k(p)}$, and its chosen period is $\log r_k(p)$.
\end{theorem}

\begin{proof}
A morphism $[L\to G]\to[L'\to G']$ of $1$-motives is a commutative
square \cite[Definition~8.1]{HuberWustholz2022}.  Here
\[
 \chi_p(u_{\mathrm{ar}}(n))
 =(2^n)^{A(p)}(3^n)^{-\ell(p)}=\alpha_p^n,
\]
so \eqref{eq:kummer-character-morphism} is such a square.  Moreover
\[
 \chi_p^*\frac{dz}{z}=A(p)\frac{dx}{x}-\ell(p)\frac{dy}{y}.
\]
Integration on $\gamma$ proves \eqref{eq:kummer-path-period}.  The prefix
statement is the same calculation with exponent vector
$(-A_k(p),k-1)$.
\end{proof}

Cartier duality makes the fixed rank-two object equally explicit.  The dual
of $M_{\mathrm{ar}}$ has the form
\begin{equation}
\label{eq:arithmetic-one-motive-dual}
 M_{\mathrm{ar}}^\vee=
 [\mathbb Z^2\longrightarrow\mathbf G_m],
 \qquad (r,s)\longmapsto2^r3^s.
\end{equation}
Thus $(A(p),-\ell(p))$ evaluates to $\alpha_p$.  Formula
\eqref{eq:arithmetic-one-motive-dual} is the specialization of the character
lattice in Cartier duality, not a second unidentified motive
\cite[Section~1.8]{BarbieriVialeKahn2016}.

\subsection{Cayley pairs and the exact Nori arrows}

Let $K=\mathbb Q(i)$ and
\[
 U_K=\mathbf P^1_K\setminus\{i,-i\}.
\]

\begin{theorem}[Cayley isomorphism of pairs]
\label{thm:cayley-pair-isomorphism}
The rational maps
\begin{equation}
\label{eq:cayley-map-and-inverse}
 \kappa(z)=\frac{1+iz}{1-iz},
 \qquad
 \kappa^{-1}(w)=\frac{w-1}{i(w+1)}
\end{equation}
are inverse isomorphisms $U_K\simeq\mathbf G_{m,K}$.  For
$r\in\mathbb Q_{>0}$ and
$\alpha(r)=(1+ir)/(1-ir)$ they give an isomorphism of pairs
\begin{equation}
\label{eq:cayley-pair-map}
 \kappa:(U_K,\{0,r\})\overset{\sim}{\longrightarrow}
 (\mathbf G_{m,K},\{1,\alpha(r)\}),
\end{equation}
and
\begin{equation}
\label{eq:cayley-differential}
 \kappa^*\frac{dw}{w}=\frac{2i\,dz}{1+z^2}.
\end{equation}
Consequently the functoriality edge in Nori's effective diagram has the
contravariant direction
\begin{equation}
\label{eq:cayley-nori-edge}
 (\mathbf G_{m,K},\{1,\alpha(r)\},1)\longrightarrow
 (U_K,\{0,r\},1).
\end{equation}
\end{theorem}

\begin{proof}
Solving $w=(1+iz)/(1-iz)$ gives the displayed inverse.  The two deleted
points map to $0$ and $\infty$, so the restriction is an isomorphism to
$\mathbf G_m$.  It sends $0$ to $1$ and $r$ to $\alpha(r)$, proving the
pair statement.  Logarithmic differentiation gives
\[
 \frac{i\,dz}{1+iz}+\frac{i\,dz}{1-iz}
 =\frac{2i\,dz}{1+z^2}.
\]
Definition~9.1.1 of \cite{HuberMullerStach2025} sends a map of pairs
contravariantly on relative cohomology, which fixes the edge direction.
\end{proof}

Over $\mathbb Q$ itself, put
\[
 U_{\mathbb Q}=\mathbf P^1_{\mathbb Q}\setminus V(z^2+1).
\]
The form $dz/(1+z^2)$ is regular on $U_{\mathbb Q}$, including at infinity,
and $0,r$ are rational points.  Hence
\[
 \int_0^r\frac{dz}{1+z^2}=\arctan r
\]
is a $1$-period of the pair $(U_{\mathbb Q},\{0,r\},1)$.
Proposition~12.10 of \cite{HuberWustholz2022} identifies this with the period
world of $1$-motives.  More explicitly, with the local coordinate $t=1/z$ at
infinity,
\[
 \frac{dz}{1+z^2}=-\frac{dt}{1+t^2},
\]
which proves regularity there.  Removing infinity is therefore unnecessary,
and doing so changes the Cayley target to $\mathbf G_m\setminus\{-1\}$.

Path concatenation by itself is not a Nori edge.  A Nori functoriality edge
requires a map of algebraic pairs, and a coboundary edge requires a nested
triple of closed subvarieties.  The exact replacement for additive path
cocycles uses multiplication.

\begin{proposition}[multiplication diagram for concatenated characters]
\label{prop:kummer-character-multiplication}
For composable paths $p,q$,
\[
 \chi_{q\circ p}=\chi_q\chi_p,
 \qquad \alpha_{q\circ p}=\alpha_q\alpha_p.
\]
Let
\[
 N_{p,q}=
 [\mathbb Z\longrightarrow\mathbf G_m^2],
 \qquad n\longmapsto(\alpha_p^n,\alpha_q^n).
\]
There is a factorization by morphisms of $1$-motives
\begin{equation}
\label{eq:kummer-multiplication-factorization}
 M_{\mathrm{ar}}
 \xrightarrow{(\mathrm{id},(\chi_p,\chi_q))}N_{p,q}
 \xrightarrow{(\mathrm{id},m)}K(\alpha_p\alpha_q),
 \qquad m(z_1,z_2)=z_1z_2.
\end{equation}
At the level of Nori pairs, the map
\[
 m:(\mathbf G_m^2,\{(1,1),(\alpha_p,\alpha_q)\})
 \longrightarrow(\mathbf G_m,\{1,\alpha_p\alpha_q\})
\]
gives the contravariant functoriality edge
\[
 (\mathbf G_m,\{1,\alpha_p\alpha_q\},1)
 \xrightarrow{m^*}
 (\mathbf G_m^2,\{(1,1),(\alpha_p,\alpha_q)\},1),
\]
and
\begin{equation}
\label{eq:dlog-multiplication}
 m^*\frac{dz}{z}
 =\operatorname{pr}_1^*\frac{dz}{z}
  +\operatorname{pr}_2^*\frac{dz}{z}.
\end{equation}
Thus this declared pair map, rather than concatenation of singular paths,
is the exact diagrammatic source of period additivity.
\end{proposition}

\begin{proof}
Lengths and total exponents add under composition, which proves the two
character identities.  Both squares in
\eqref{eq:kummer-multiplication-factorization} commute on $n\in\mathbb Z$:
the first maps $u_{\mathrm{ar}}(n)$ to
$(\alpha_p^n,\alpha_q^n)$, and multiplication maps that point to
$(\alpha_p\alpha_q)^n$.  The displayed map sends the two marked points into
the two marked points, so it is a map of pairs.  Logarithmic differentiation
proves \eqref{eq:dlog-multiplication}.
\end{proof}

For $r>0$,
\[
 \operatorname{Im}\alpha(r)=\frac{2r}{1+r^2}>0.
\]
It therefore cannot equal the positive rational number
$2^{A(p)}3^{-\ell(p)}$.  The Cayley pair for a prefix angle and the rational
Kummer pair for the path weight are linked by the toric constructions above,
but they are not identified by an unproved equality of marked points.

\begin{nonclaim}
The section constructs the arithmetic groupoid, the germ groupoid with
compact unit space,
their faithful cocycle-preserving functor, the torus isomorphism, the Kummer
morphisms, the Cayley pair isomorphism, and the multiplication edges.  It
does not construct a universal path-indexed Nori category, a classification
of equilibrium states, a motivic Galois action on the path algebra, or a
Collatz termination theorem.
\end{nonclaim}
